\documentclass[a4paper,11pt]{article}

\usepackage[utf8]{inputenc}
\usepackage[T1]{fontenc}
\usepackage{lmodern}
\usepackage{amsmath,amssymb,amsthm}
\usepackage{mathtools}
\usepackage{hyperref}
\usepackage[margin=2.5cm]{geometry}
\usepackage{enumitem}

\theoremstyle{definition}
\newtheorem{definition}{Definition}[section]
\newtheorem{theorem}{Theorem}[section]
\newtheorem{remark}{Remark}[section]

\title{\textbf{CertiQ Dispatcher}\\z3 Formal Definition}
\author{CertiQ Project}
\date{\today}

\begin{document}

\maketitle

\tableofcontents
\newpage

% ============================================================
\section*{Preface}
\addcontentsline{toc}{section}{Preface}
% ============================================================

This document defines the z3 CertiQ Dispatcher architecture.

The purpose of z3 is not to patch the old model names. It is to state the
architecture from first principles before any implementation rewrite happens.
The core object is a certified online dispatcher for assigning arriving work to
heterogeneous capacity-constrained resources.

\subsection*{z3 Thesis}

CertiQ Dispatcher is a certified differentiable architecture for online
assignment:

\begin{equation}\label{eq:thesis}
\pi^\Theta(Q,\mu,\xi)
=
\mathcal C_Q\bigl(
\mathcal A_\Theta(Q,\mu,\xi),
\mathcal B(Q,\mu)
\bigr).
\end{equation}

Here:
\begin{enumerate}[label=(\roman*)]
  \item \(\mathcal B\) is the certified base dispatch geometry;
  \item \(\mathcal A_\Theta\) is the learned assignment proposal;
  \item \(\mathcal C_Q\) is the state-dependent certificate operator;
  \item \(\pi^\Theta\) is the final dispatch distribution.
\end{enumerate}

The architecture is not defined by preview implementation variants. It is the
composition of base geometry, learned proposal, and certificate operator.

\subsection*{Relationship to z2}

The z2 material remains the current proof package for the CTMC model:
\begin{itemize}
  \item \texttt{docs/z2/formal\_math/01\_backbone\_stability\_and\_constant.md}
  \item \texttt{docs/z2/formal\_math/02\_gate\_inheritance\_theorems.md}
\end{itemize}
z3 reuses proven z2 facts where they apply, but it changes the architecture
language. In z3, certificate mechanisms are first-class operators, not gates
attached after a neural policy has already been designed.

\subsection*{Non-Goals}

This document does not implement the rewrite. It also does not prove new
stability theorems beyond the z2 proof package. New claims introduced by z3 must
remain theorem obligations until separately proved.

% ============================================================
\section{Architecture Principles}
% ============================================================

\subsection{Design Target}

CertiQ Dispatcher is designed for one primitive:
\begin{quote}
assign one arriving unit of work to one resource, online, under heterogeneous
service capacity, while preserving an auditable stability certificate.
\end{quote}

The architecture should be judged by whether each component helps this
primitive. A component is valid only if it has a clear role in at least one of:
\begin{enumerate}
  \item representing resources without depending on arbitrary labels;
  \item comparing current workload against capacity;
  \item proposing a performance-improving assignment;
  \item enforcing a certificate before dispatch;
  \item exposing diagnostics that allow the certificate to be audited.
\end{enumerate}

\subsection{Architectural Discipline}

The lesson to borrow from the Transformer is not the attention mechanism
itself. The lesson is architectural compression: define a small number of
operations that directly match the task primitive, then reuse those operations
consistently.

For CertiQ Dispatcher, the primitive is not token-to-token contextualization.
It is capacity-aware assignment. Therefore the central operation is:
\[
\text{certified assignment}
=
\text{proposal}
\quad\text{filtered through}\quad
\text{state-dependent admissibility}.
\]

Every z3 block must answer:
\begin{enumerate}
  \item What mathematical object does this block produce?
  \item What invariance or constraint does it preserve?
  \item What failure mode does it prevent?
  \item What diagnostic can verify that it behaved correctly?
\end{enumerate}

\subsection{Three-Layer Architecture}

The architecture has three layers.

\subsubsection{Geometry Layer}

The geometry layer maps queue state and capacity into a stable dispatch energy.
It defines the coordinates in which backlog is compared:
\[
y_i(Q)=Q_i/\mu_i^\beta.
\]
It also defines the certified base policy \(p^{\mathrm{cert}}\). This policy
must have a theorem-facing arrival envelope:
\[
A_{p^{\mathrm{cert}}}(Q)\le m(Q)+C.
\]

\subsubsection{Proposal Layer}

The proposal layer learns performance corrections. It may use context
\(\xi\), attention, pooling, learned costs, compatibility features, or other
domain signals.

The proposal layer is not the source of certification. It is the source of
adaptivity and performance.

\subsubsection{Certificate Layer}

The certificate layer maps an unconstrained or partially constrained proposal
into a certified final distribution.

The certificate layer is where the architecture earns the name CertiQ
Dispatcher. If the final policy violates the certificate, the architecture has
failed, even if training loss improves.

\subsection{Required Invariances}

Resource labels have no intrinsic meaning. The architecture must be
permutation equivariant:
\[
\pi(\sigma Q,\sigma\mu,\sigma\xi)=\sigma\pi(Q,\mu,\xi).
\]

This requirement affects all blocks:
\begin{enumerate}
  \item local feature maps must be shared across resources;
  \item global summaries must be permutation invariant;
  \item resource scores must be produced equivariantly;
  \item the certificate operator must commute with resource permutation;
  \item diagnostics must reorder consistently with the state.
\end{enumerate}

\subsection{First-Class Certificate Boundary}

z3 forbids treating certification as a training regularizer.

The final action distribution must satisfy the certificate at inference time by
construction. A drift penalty may help training, but it cannot substitute for
projection, fallback, or another exact certificate operator.

\subsection{What Makes The Design Elegant}

An elegant CertiQ Dispatcher architecture should have:
\begin{enumerate}
  \item one canonical forward equation;
  \item one resource-set interface;
  \item one diagnostics object;
  \item pluggable proposal encoders;
  \item pluggable certificate operators;
  \item theorem claims stated only at the certificate boundary.
\end{enumerate}
The implementation may expose multiple model presets, but those presets must
be parameterizations of the same architecture.

% ============================================================
\section{Formal Dispatch Model}
% ============================================================

\subsection{Resources}

There are \(N\ge 1\) resources. Resource \(i\) has service capacity
\(\mu_i>0\). The capacity vector is:
\[
\mu=(\mu_1,\ldots,\mu_N).
\]
Total service capacity is:
\[
\Lambda=\sum_{i=1}^N\mu_i.
\]
Each resource may carry context:
\[
\xi_i\in\mathbb R^{d_\xi}.
\]
The context may describe metadata, cost, compatibility, location, energy,
hardware, or other domain information. Context does not automatically enter the
certificate. It enters certification only through quantities explicitly named
in a theorem.

\subsection{Arrivals And State}

Jobs arrive online. In the primary certifiable model, arrivals are Poisson with
rate:
\[
\lambda>0.
\]
The state is:
\[
Q=(Q_1,\ldots,Q_N)\in\mathbb Z_+^N,
\]
where \(Q_i\) is the number of active or queued jobs at resource \(i\).

The primary CTMC certificate assumes:
\[
\lambda<\Lambda.
\]
This is the subcritical load condition. It is necessary for the aggregate
system but does not certify arbitrary policies.

\subsection{Actions}

At each arrival, the dispatcher observes:
\[
(Q,\mu,\xi)
\]
and outputs:
\[
\pi(Q,\mu,\xi)\in\Delta_N,
\]
where:
\[
\Delta_N=\{p\in\mathbb R_+^N:\textstyle\sum_i p_i=1\}.
\]
The arriving job is assigned to resource \(i\) with probability \(\pi_i\).

\subsection{Service Model}

For the primary certifiable CTMC model, service completions are exponential.
When \(Q_i>0\), resource \(i\) completes one job at rate \(\mu_i\).

The generator under policy \(\pi\) is:
\[
(\mathcal L_\pi f)(Q)
=
\lambda\sum_i\pi_i(Q,\mu,\xi)[f(Q+e_i)-f(Q)]
+
\sum_i\mu_i\mathbf 1_{\{Q_i>0\}}[f(Q-e_i)-f(Q)].
\]
The policy controls only the arrival term. Service dynamics are fixed by
\(\mu\).

\subsection{Objective}

The default performance objective is average queueing cost:
\[
J(\pi)
=
\limsup_{T\to\infty}
\frac1T
\mathbb E_\pi
\left[
\int_0^T c(Q(t),\mu,\xi)\,dt
\right].
\]
The canonical cost is total backlog:
\[
c(Q)=\sum_i Q_i.
\]
Weighted costs are allowed when weights are positive and bounded:
\[
c(Q,\mu,\xi)=\sum_i w_i(\mu,\xi)Q_i.
\]
Certification is about stability and admissibility. It is not a proof of
optimality.

\subsection{Adapter Contract}

A domain adapter maps a real system into:
\[
\mathcal D=(N,\lambda,\mu,\xi,Q,c).
\]
An adapter must state:
\begin{enumerate}
  \item what counts as an arrival;
  \item what counts as a resource;
  \item how \(\mu_i\) is measured;
  \item what \(Q_i\) counts;
  \item what service law is assumed;
  \item what context \(\xi_i\) means;
  \item what cost is optimized;
  \item whether the CTMC certificate applies exactly.
\end{enumerate}
If the adapter violates the CTMC assumptions, CertiQ Dispatcher may still be
used empirically, but the CTMC theorem must not be claimed.

% ============================================================
\section{CertiQ Dispatcher Architecture}
% ============================================================

\subsection{Canonical Equation}

The z3 architecture is:
\[
\pi^\Theta(Q,\mu,\xi)
=
\mathcal C_Q\bigl(
\mathcal A_\Theta(Q,\mu,\xi),
\mathcal B(Q,\mu)
\bigr).
\]
This equation defines the architecture.

\(\mathcal B\) is the base certified geometry. \(\mathcal A_\Theta\) is the
learned assignment proposal. \(\mathcal C_Q\) is the state-dependent
certificate operator. The final policy \(\pi^\Theta\) is the only distribution
used for dispatch.

\subsection{Resource Coordinates}

The architecture begins by computing capacity-normalized workload:
\[
y_i(Q)=Q_i/\mu_i^\beta,
\qquad
m(Q)=\min_i y_i(Q).
\]
These coordinates are not optional metadata. They are the Lyapunov geometry
used by the certificate layer.

\subsection{Base Certified Geometry}

The base geometry produces logits:
\[
u_i^{\mathrm{cert}}
=
\gamma\log\mu_i
-
\alpha\frac{Q_i+c}{\mu_i^\beta}.
\]
The base policy is:
\[
p^{\mathrm{cert}}
=
\operatorname{softmax}(u^{\mathrm{cert}}).
\]
For the z2 CTMC theorem, this is the analytic backbone and satisfies:
\[
A_{p^{\mathrm{cert}}}(Q)\le m(Q)+C.
\]
In z3 language, the important object is not the old name ``backbone''. The
important object is a certified base distribution with an arrival-envelope
constant.

\subsection{Learned Assignment Proposal}

The proposal layer builds an equivariant score correction from resource-local
and global information.

Local features:
\[
s_i=[
\log(1+Q_i),
\log\mu_i,
y_i,
\mu_i/\Lambda,
\xi_i
].
\]

Shared local encoding:
\[
z_i^0=f_{\mathrm{loc}}(s_i).
\]

Invariant global context:
\[
g=f_{\mathrm{glob}}(\{z_i^0\}_{i=1}^N).
\]

Equivariant resource update:
\[
z_i^1=f_{\mathrm{res}}([z_i^0,g]).
\]

Proposal correction:
\[
r_i^\Theta=R_{\max}\tanh(h(z_i^1)).
\]

Proposal policy:
\[
p^\Theta_{\mathrm{prop}}
=
\operatorname{softmax}(u^{\mathrm{cert}}+r^\Theta).
\]

The proposal is anchored to the certified base logits. It is not a free
black-box dispatcher.

\subsection{Raw Preference For Learning}

The proposal layer may also produce a raw usage preference:
\[
\eta_{\mathrm{raw}}\in[0,1].
\]
\(\eta_{\mathrm{raw}}\) is not a certificate. It is only the learned desire to
use the proposal.

\subsection{Candidate Mixture}

The candidate mixture is:
\[
\pi_\eta
=
(1-\eta)p^{\mathrm{cert}}
+
\eta p^\Theta_{\mathrm{prop}}.
\]
The certificate operator chooses or modifies \(\eta\), or more generally
projects the candidate policy, so the final dispatch distribution is
admissible.

\subsection{Certificate Operator}

The certificate operator receives:
\begin{enumerate}
  \item current state \(Q\);
  \item base distribution \(p^{\mathrm{cert}}\);
  \item proposal distribution \(p^\Theta_{\mathrm{prop}}\);
  \item raw usage preference \(\eta_{\mathrm{raw}}\);
  \item certificate constants and diagnostics.
\end{enumerate}

It returns:
\begin{enumerate}
  \item final policy \(\pi^\Theta\);
  \item final gate or projection variables;
  \item certificate diagnostics.
\end{enumerate}

The final policy must be the distribution used by the simulator, trainer,
evaluator, and deployment path.

\subsection{Architecture Presets}

z3 may expose presets:
\begin{enumerate}
  \item \textbf{CertiQ Dispatcher-F}: tail fallback certificate operator;
  \item \textbf{CertiQ Dispatcher-P}: drift-envelope projection operator;
  \item \textbf{CertiQ Dispatcher-X}: uncertified ablation.
\end{enumerate}
These are presets of one architecture, not separate architectures.

The certified presets must share the same forward interface and diagnostics
contract.

% ============================================================
\section{Certificate Layer}
% ============================================================

\subsection{Certificate Quantity}

For any policy \(\pi\), define:
\[
A_\pi(Q)=\sum_i \pi_i(Q,\mu,\xi)\,y_i(Q).
\]

The certified envelope is:
\[
B(Q)=m(Q)+C.
\]

The certificate condition is:
\[
A_{\pi^\Theta}(Q)\le B(Q).
\]

For the current z2 CTMC theorem, \(C=C_B\), the exact constant from the
analytic backbone proof.

\subsection{Admissible Dispatch Set}

At state \(Q\), define the admissible simplex:
\[
\mathcal P_{\mathrm{cert}}(Q)
=
\{\pi\in\Delta_N:A_\pi(Q)\le B(Q)\}.
\]

The certificate operator must return:
\[
\pi^\Theta(Q,\mu,\xi)\in\mathcal P_{\mathrm{cert}}(Q).
\]

This makes certification an action-space constraint, not an auxiliary loss.

\subsection{Tail Fallback Operator}

Define:
\[
S_\beta(Q)=\sum_i Q_i/\mu_i^\beta.
\]

For finite \(R_{\mathrm{cert}}\), the fallback operator sets:
\[
\eta=
\eta_{\mathrm{raw}}\mathbf 1_{\{S_\beta(Q)\le R_{\mathrm{cert}}\}}.
\]

Then:
\[
\pi^\Theta
=
(1-\eta)p^{\mathrm{cert}}+\eta p^\Theta_{\mathrm{prop}}.
\]

When \(S_\beta(Q)>R_{\mathrm{cert}}\), the policy equals the certified base
exactly.

This operator is conservative. It certifies by making the learned proposal a
finite-region perturbation of the base policy.

\subsection{Drift-Envelope Projection Operator}

Compute:
\[
A_{\mathrm{cert}}(Q)=A_{p^{\mathrm{cert}}}(Q),
\qquad
A_{\mathrm{prop}}(Q)=A_{p^\Theta_{\mathrm{prop}}}(Q).
\]

If \(A_{\mathrm{prop}}(Q)\le A_{\mathrm{cert}}(Q)\), the proposal is no worse
than the base under the envelope coordinate and may be fully allowed:
\[
\eta_{\mathrm{safe}}=1.
\]

Otherwise:
\[
\eta_{\mathrm{safe}}
=
\left[
\frac{B(Q)-A_{\mathrm{cert}}(Q)}
{A_{\mathrm{prop}}(Q)-A_{\mathrm{cert}}(Q)}
\right]_0^1.
\]

The final usage is:
\[
\eta=\min(\eta_{\mathrm{raw}},\eta_{\mathrm{safe}}).
\]

The final policy:
\[
\pi^\Theta=(1-\eta)p^{\mathrm{cert}}+\eta p^\Theta_{\mathrm{prop}}
\]
satisfies:
\[
A_{\pi^\Theta}(Q)\le B(Q).
\]

This operator certifies by pointwise envelope preservation.

\subsection{Diagnostics Contract}

Every certified forward pass must emit:
\begin{enumerate}
  \item \(A_{\mathrm{cert}}\);
  \item \(A_{\mathrm{prop}}\);
  \item \(A_{\pi^\Theta}\);
  \item \(m(Q)\);
  \item \(B(Q)\);
  \item \(B(Q)-A_{\pi^\Theta}\);
  \item raw proposal usage;
  \item final proposal usage;
  \item projection cap when applicable;
  \item fallback indicator when applicable;
  \item residual magnitude;
  \item final policy entropy;
  \item selected resource or action distribution summary.
\end{enumerate}
These are architecture diagnostics. They are not optional experiment metadata.

\subsection{Invalid Certificate Patterns}

The following do not certify the architecture:
\begin{enumerate}
  \item adding a drift penalty without exact enforcement;
  \item using stale probabilities for projection;
  \item projecting one policy and dispatching another;
  \item using smooth fallback at inference without a separate proof;
  \item claiming CTMC stability for adapters that violate the CTMC model;
  \item allowing approximate projection without a numerical tolerance contract
        and proof that the resulting error preserves drift.
\end{enumerate}

% ============================================================
\section{Training And Objectives}
% ============================================================

\subsection{Training Is Not Certification}

Training selects useful parameters. The certificate layer defines admissible
actions. These roles must remain separate.

The training objective may improve cost, imitation quality, exploration,
entropy, and proposal usefulness. It must not be used as the basis for a
stability claim unless the final policy is also certified by construction.

\subsection{Objective Hierarchy}

The z3 training objective should be organized by priority:
\begin{enumerate}
  \item maintain exact certificate enforcement in the forward path;
  \item reduce average queueing cost;
  \item learn useful proposal corrections over the certified base geometry;
  \item preserve resource permutation equivariance;
  \item keep diagnostics auditable and non-silent.
\end{enumerate}
This hierarchy prevents the learned proposal from becoming a hidden replacement
for the certified architecture.

\subsection{Suggested Loss Components}

The training loss may include:
\[
\mathcal L
=
\omega_{\mathrm{roll}}\mathcal L_{\mathrm{roll}}
+
\omega_{\mathrm{bc}}\mathcal L_{\mathrm{bc}}
+
\omega_{\mathrm{res}}\mathcal L_{\mathrm{res}}
+
\omega_{\mathrm{usage}}\mathcal L_{\mathrm{usage}}
+
\omega_{\mathrm{ent}}\mathcal L_{\mathrm{ent}}
+
\omega_{\mathrm{cert}}\mathcal L_{\mathrm{cert}}.
\]

The certificate penalty:
\[
\mathcal L_{\mathrm{cert}}
=
\mathbb E\bigl[(A_{\pi^\Theta}(Q)-B(Q))_+^2\bigr]
\]
should be zero when exact projection is active. If it is not zero, the
implementation is wrong or the numerical tolerance is incorrectly declared.

\subsection{Proposal Learning}

The learned proposal should be trained as a correction over the certified base,
not as an independent dispatcher.

Acceptable proposal targets include:
\begin{enumerate}
  \item rollout-improved certified policy;
  \item dynamic-programming policy on tiny systems;
  \item audited queueing heuristic;
  \item cost-improving residual discovered by on-policy learning.
\end{enumerate}
The target must be labeled. A heuristic target is not a theorem.

\subsection{Evaluation}

Evaluation must report both performance and certificate metrics.

Performance metrics:
\begin{enumerate}
  \item average queue length;
  \item average cost;
  \item tail backlog quantiles;
  \item maximum observed backlog;
  \item latency proxy where applicable.
\end{enumerate}

Certificate metrics:
\begin{enumerate}
  \item violation rate;
  \item minimum drift slack;
  \item average drift slack;
  \item projection activation rate;
  \item fallback activation rate;
  \item proposal usage;
  \item residual magnitude.
\end{enumerate}

No result should be presented as certified without certificate metrics.

\subsection{Ablations}

Ablations should test the architecture thesis:
\begin{enumerate}
  \item base policy only;
  \item proposal without certificate;
  \item certificate with weak proposal;
  \item projection versus fallback;
  \item equivariant proposal versus order-dependent proposal;
  \item context-free versus context-aware proposal.
\end{enumerate}
The uncertified ablation must be labeled as uncertified.

\subsection{Curriculum}

Recommended curriculum:
\begin{enumerate}
  \item verify base certificate quantities;
  \item train proposal to imitate base or audited heuristic;
  \item enable exact certificate operator during all policy rollouts;
  \item optimize rollout cost;
  \item audit state banks before reporting results;
  \item only then expand to approximate adapters or richer domains.
\end{enumerate}

% ============================================================
\begin{thebibliography}{9}
% ============================================================

\bibitem{vaswani2017attention}
A.~Vaswani, N.~Shazeer, N.~Parmar, J.~Uszkoreit, L.~Jones, A.~N.~Gomez,
\L.~Kaiser, and I.~Polosukhin.
\newblock Attention is all you need.
\newblock \textit{arXiv preprint arXiv:1706.03762}, 2017.

\bibitem{zaheer2017deep}
M.~Zaheer, S.~Kottur, S.~Ravanbakhsh, B.~Poczos, R.~Salakhutdinov,
and A.~Smola.
\newblock Deep sets.
\newblock In \textit{NeurIPS}, 2017.

\bibitem{lee2019set}
J.~Lee, Y.~Lee, J.~Kim, A.~Kosiorek, S.~Choi, and Y.~W.~Teh.
\newblock Set transformer: A framework for attention-based
permutation-invariant neural networks.
\newblock In \textit{ICML}, 2019.

\bibitem{meyn2012markov}
S.~P.~Meyn and R.~L.~Tweedie.
\newblock \textit{Markov Chains and Stochastic Stability}.
\newblock Cambridge University Press, 2nd~edition, 2009.

\bibitem{tassiulas1992stability}
L.~Tassiulas and A.~Ephremides.
\newblock Stability properties of constrained queueing systems and scheduling
policies for maximum throughput in multihop radio networks.
\newblock \textit{IEEE Trans.~Automatic Control}, 37(12):1936--1948, 1992.

\bibitem{certiqz2}
CertiQ Project.
\newblock z2 proof package.
\newblock \path{docs/z2/formal_math/01_backbone_stability_and_constant.md},
\path{docs/z2/formal_math/02_gate_inheritance_theorems.md}.

\end{thebibliography}

\end{document}
